\Chapter{Testowanie i weryfikacja}
\label{ch:verification}

\section{Co sprawdzono}

Testy działają bez internetu dzięki \code{pytest-socket}; połączenia
internetowe są blokowane, a lokalne gniazda uniksowe pozostają dostępne.
Fixtures to niewielki, wersjonowany materiał wejściowy do testów.
Zawierają pełne surowe tabele Bahrajnu 2022 (1125 wierszy) i Arabii
Saudyjskiej 2022 (820 wierszy), a nie losowo wymyślone próbki zastępujące
rzeczywiste dane. Oprócz nich testy budują małe przypadki syntetyczne,
pozwalające celowo uszkodzić pojedynczy warunek.

\begin{table}[htbp]
\centering\small
\caption{Zakres istniejących testów i granice ich interpretacji.}
\label{tab:tests}
\begin{tabularx}{\textwidth}{L{0.25\textwidth}Y Y}
\toprule
Obszar & Sprawdzana własność & Czego wynik nie potwierdza \\
\midrule
Kontrakty & Typy, klucze, wartości, zgodność \code{row_id}, kształty i dopełnienie tablic. & Przyczynowości przyszłego generatora cech. \\
Fixtures i adapter & Odczyt Parquet, zachowanie braków, korekta tekstowych nulli, historyczne numery kierowców. & Kompletności historii opon w pełnym zbiorze. \\
Snapshot & Pakowanie i odczyt bez cache, niezmienność nazw, reakcja na brak źródeł. & Akceptacji danych do uczenia. \\
Weryfikator & Zmienione/brakujące pliki, niebezpieczne ścieżki, strumienie i semantyka metadanych. & Prawdziwości każdej wizyty jako wymiany opon. \\
CLI i konfiguracje & Lista poleceń, kody zakończenia, odczyt YAML, jawne ograniczenie sieci. & Działania niewdrożonych poleceń modelowych. \\
\bottomrule
\end{tabularx}
\end{table}

Zachowany wynik końcowych kontroli autora T2 na rewizji
\code{e4d484d} podaje 209 poprawnych testów oraz poprawne kontrole Ruff.
Nie jest to nowy pomiar wykonany podczas pisania raportu. Dowód merge
wiąże tę rewizję z \code{5ba376f}; CI dla kodu przed scaleniem zakończyło
się powodzeniem~\cite{t2ci,t2merge}. Weryfikacja świeżo rozpakowanego
archiwum przez właściwe CLI również zakończyła się kodem 0, przy blokadzie
FastF1, cache, uwierzytelniania i sieci. Kopie testowe bez zadeklarowanego
strumienia lub z całkowicie pustą treścią \code{Status}/\code{Message}
zostały odrzucone kodem 1 mimo zgodnych, przeliczonych skrótów.

Kandydat snapshotu powstał wcześniej na rewizji \code{e96ae05}.
Nowsza poprawka weryfikatora go nie wytworzyła ponownie, tylko sprawdziła.
Rewizja producenta danych i rewizja programu weryfikującego są więc
różnymi elementami pochodzenia. Ich pełne identyfikatory i skróty
artefaktów znajdują się w mapie dowodów.

\section{Wnioski z kontroli FastF1}

Audyt zachował związek \code{PitInTime} z okrążeniem wjazdowym
i \code{PitOutTime} z wyjazdowym. Nie stwierdza to jednak, że każdy
znacznik czasu musi leżeć w granicach zapisanej długości okrążenia.
Odnotowano 97 znaczników wjazdu późniejszych niż odpowiadający im koniec
okrążenia. Nie przeniesiono ich arbitralnie do następnego wiersza.
Reguła przypisywania zdarzeń wymaga zachowania źródłowego kontekstu.

\code{IsAccurate} nie może służyć do wyboru wierszy z etykietą.
W badanych danych okrążenia wjazdowe i wyjazdowe nie są dokładne według
tej flagi, więc ich usunięcie zmieniłoby populację zdarzeń. Flaga ma
maskować tylko cechy czasowe, z przyczynowym przeniesieniem wcześniejszych
wartości w obrębie wyścigu i kierowcy. To reguła dla przyszłej implementacji,
nie opis wykonanego przetwarzania.

Istotne jest także dokładne znaczenie statusu SC. Po poprzednim statusie
równym dokładnie \code{4} wszystkie 2806 obserwowanych wierszy było
niedokładnych. Po dowolnym statusie zawierającym \code{4}, również
wielokodowym, 1017 z 5260 wierszy miało \code{IsAccurate=True}.
Nie wolno uogólnić reguły dla statusu dokładnego na każdy status mieszany.
Jest to ustalenie audytu FastF1 3.9.0a0, wymagające doprecyzowania
starszego opisu reguły~\cite{findings}.

\section{Ochrona przed wyciekiem: stan i plan}

Wyciekiem informacji nazywamy wykorzystanie przy prognozie danych, które
w rzeczywistym momencie jej sporządzania nie były dostępne, albo
dopasowanie przetwarzania na wydzielonych danych oceny.
Obecnie działają mechanizmy zachowania surowych obserwacji, znaczników
czasu, pochodzenia i spójności formatu. Nie ma jeszcze generatora cech
ani testu obcięcia całego zbioru do $T$. Nagłówek reguł w dokumencie
decyzji nie stanowi dowodu, że wszystkie przyszłe testy już istnieją.

Docelowa kontrola ma ponownie obliczać cechy po usunięciu obserwacji
po $T$ ze wszystkich strumieni i od wszystkich kierowców. Wynik musi
odpowiadać obliczeniu na pełnym wejściu. Osobne testy mają sprawdzać
przesunięcie etykiety na $t+1$, obojętność na przyszły wjazd rywala,
nieprzekraczanie granic kierowcy i wyścigu w oknie oraz dopasowanie
standaryzacji wyłącznie na części uczącej. Wierszy nie wolno wybierać
za pomocą \code{IsAccurate} ani usuwać na podstawie braku czasu
okrążenia~\cite{decisions}.

\section{Protokół przyszłych eksperymentów}

Przyjęty podział jest chronologiczny: train obejmuje sezony 2022--2024,
val rundy 1--12 sezonu 2025, a test rundy 13--24 tego sezonu. Manifest
podziału ma powstać przed wykluczeniami. Nie przewidziano losowego
podziału wierszy ani ponownego uczenia na train+val.
Statystyki standaryzacji mają pochodzić tylko z train i pozostawać
zamrożone dla val, test oraz późniejszego wnioskowania.

Konfiguracje zakładają porównanie stałej częstości klasy z train,
regresji logistycznej i LightGBM z perceptronem wielowarstwowym (MLP),
siecią rekurencyjną z bramkowaną komórką GRU i czasową siecią
konwolucyjną (TCN).
MLP ma przetwarzać pojedynczy wektor cech, a GRU i przyczynowa TCN
okno historii. To plany architektury; nie istnieje jeszcze ich
\code{forward}, pętla uczenia ani zapis wag. Sieci mają zwracać jeden
logit, czyli wynik przed funkcją sigmoidalną zamieniającą go na
prawdopodobieństwo. Referencyjny trening ma używać AdamW i
\code{BCEWithLogitsLoss} bez zwiększania wagi klasy dodatniej.

Podstawową metryką wyboru modelu będzie Average Precision (AP),
podawana razem z częstością klasy dodatniej. AP ocenia ranking wyników
przy zmianie progu decyzji. Próg alarmu ma maksymalizować F1 na val,
z wyborem większego progu przy remisie, a następnie zostać zamrożony.
F1 łączy precyzję, czyli udział trafnych alarmów, i czułość, czyli
udział wykrytych zdarzeń. Dodatkowa ocena zdarzeniowa ma dopuszczać
wcześniejszy alarm, ale nie alarm spóźniony względem okna etykiety.
Niepewność ma być oceniana przez wielokrotne losowanie ze zwracaniem
całych wyścigów (bootstrap), co zachowuje zależności między ich
okrążeniami~\cite{decisions,contracts}.

Nie przeprowadzono uczenia, porównania modeli ani predykcyjnej oceny
val/test. Zablokowany test wymaga osobnej zgody człowieka po zamrożeniu
konfiguracji. Również ewentualne AP lub F1 powyżej 0.6 na val/test wymaga
audytu wycieku przed publikacją wyników. Analiza jakości danych z tych
sezonów nie jest predykcyjną oceną na zbiorze testowym.
